\SourceLeaf{016}
\[
\begin{aligned}
F\mu&=\lambda m=\psi\mu\\[3pt]
F'\mu&=\frac{\lambda m}{(1-f)\mu}=\psi'\mu
\end{aligned}
\]
Quare quum \(\psi'y\), aucta ipsa \(y\), continuo crescat (saltem non decrescat, quod semper subintelligendum), \(F'y\) vero constans sit, differentia \(\psi'y-F'y=\frac{\mathrm{d}(\psi y-Fy)}{\mathrm{d}y}\) erit positiva pro valoribus ipsius \(y\) maioribus quam \(\mu\), negativa pro minoribus. Hinc facile colligitur, \(\psi y-Fy\) semper esse quantitatem positivam, adeoque \(\psi y\) semper erit absolute maior, saltem non minor, quam \(Fy\), certe quamdiu valor ipsius \(Fy\) est positivus, i. e. ab \(y=\mu f\) usque ad \(y=1\). Hinc valor integralis \(\int(Fy)^2\mathrm{d}y\) ab \(y=\mu f\) usque ad \(y=1\) erit minor valore integralis \(\int(\psi y)^2\mathrm{d}y\) inter eosdem limites, adeoque a potiori minor valore huius integralis ab \(y=0\) usque ad \(y=1\), qui fit \(=mm\). At valor integralis prioris invenitur
\[
=\frac{\lambda\lambda mm(1-\mu f)^3}{3\mu\mu(1-f)^2}
\]
unde colligimus, \(\lambda\lambda\) esse minorem quam \(\frac{3\mu\mu(1-f)^2}{(1-\mu f)^3}\), ubi \(f\) est quantitas inter \(0\) et \(1\) iacens. Iam valor fractionis \(\frac{3\mu\mu(1-f)^2}{(1-\mu f)^3}\), cuius differentiale, si \(f\) tamquam quantitas variabilis consideratur, fit \(=\)
\[
-\frac{3\mu\mu(1-f)}{(1-\mu f)^4}.(2-3\mu+\mu f)\mathrm{d}f
\]
continuo decrescit, dum \(f\) a valore \(0\) usque ad valorem \(1\) transit, quoties \(\mu\) minor est quam \(\tfrac23\), adeoque valor maximus possibilis erit is, qui valori \(f=0\) respondet, puta \(=3\mu\mu\), ita ut in hoc casu \(\lambda\) certo fiat minor vel non maior quam \(\mu\BareRoot{3}\). Q.~E.~P. Contra quoties \(\mu\) maior est quam \(\tfrac23\), valor istius fractionis erit maximus pro \(2-3\mu+\mu f=0\), i. e. pro \(f=3-\tfrac2\mu\), unde ille fit \(=\frac4{9(1-\mu)}\), adeoque in hoc casu \(\lambda\) non maior quam \(\frac2{3\BareRoot{(1-\mu)}}\). Q~E.~S.

Ita e. g. pro \(\mu=\tfrac12\) certo \(\lambda\) nequit esse maior quam \(\BareRoot{\tfrac34}\), i. e. error probabilis superare nequit limitem \(0{,}8660254m\), cui in exemplo primo art. 9 aequalis inventus est. Porro facile e theoremate nostro concluditur, \(\mu\) non esse minorem quam \(\lambda\BareRoot{\tfrac13}\), quamdiu \(\lambda\) minor sit quam \(\BareRoot{\tfrac43}\), contra \(\mu\) non esse minorem quam \(1-\frac4{9\lambda\lambda}\), pro valore ipsius \(\lambda\) maiore quam \(\BareRoot{\tfrac43}\).

\Article{11}
Quum plura problemata infra tractanda etiam cum valore integralis \(\int x^4\varphi x.\mathrm{d}x\) nexa sint, operae pretium erit, eum pro quibusdam casibus speciali\SourceLeaf{017}bus evolvere. Denotabimus valorem huius integralis ab \(x=-\infty\) usque ad \(x=+\infty\) extensi per \(n^4\).

I. Pro \(\varphi x=\tfrac1{2a}\), quatenus \(x\) inter \(-a\) et \(+a\) continetur, habemus \(n^4=\tfrac15a^4=\tfrac95m^4\).

II. In casu secundo art. 6, ubi \(\varphi x=\frac{a\mp x}{aa}\), pro valoribus ipsius \(x\) inter \(0\) et \(\pm a\), fit \(n^4=\tfrac1{15}a^4=\tfrac{12}{5}m^4\).

III. In casu tertio, ubi
\[
\varphi x=\frac{e^{-\frac{xx}{hh}}}{a\BareRoot{\pi}}
\]
invenitur per ea, quae in commentatione supra citata exponuntur, \(n^4=\tfrac34h^4=3m^4\).

Ceterum demonstrari potest, valorem ipsius \(\frac{n^4}{m^4}\) certo non esse minorem quam \(\tfrac95\), si modo suppositio art. praec. locum habeat.

\Article{12}
Denotantibus \(x,x',x''\) etc. indefinite errores observationum eiusdem generis ab invicem independentes, quorum probabilitates relativas exprimit praefixa characteristica \(\varphi\); nec non \(y\) functionem datam rationalem indeterminatarum \(x,x',x''\) etc.: integrale multiplex (I)
\[
\int\varphi x.\varphi x'.\varphi x''\mathord{.}\,\mathord{.}\,\mathord{.}\,\mathord{.}\,\mathord{.}\,\mathrm{d}x.\mathrm{d}x'.\mathrm{d}x''\mathord{.}\,\mathord{.}\,\mathord{.}\,\mathord{.}
\]
extensum per omnes valores indeterminatarum \(x,x',x''\), pro quibus valor ipsius \(y\) cadit intra limites datos \(0\) et \(\eta\), exprimet probabilitatem valoris ipsius \(y\) indefinite intra \(0\) et \(\eta\) siti. Manifesto hoc integrale erit functio ipsius \(\eta\), cuius differentiale statuemus \(=\psi\eta.\mathrm{d}\eta\), ita ut integrale ipsum fiat aequale integrali \(\int\psi\eta.\mathrm{d}\eta\) ab \(\eta=0\) incepto. Hoc pacto simul characteristica \(\psi\eta\) probabilitatem relativam cuiusvis valoris ipsius \(y\) exprimere censenda est. Quum \(x\) considerari possit tamquam functio indeterminatarum \(y,x',x''\) etc., quam statuemus
\[
=f(y,x',x''\mathord{.}\,\mathord{.}\,\mathord{.}\,\mathord{.})
\]
integrale (I) fiet
\[
\begin{aligned}
=\int\varphi.f(y,x',x''\mathord{.}\,\mathord{.}\,\mathord{.}).\frac{\mathrm{d}f(y,x',x''\mathord{.}\,\mathord{.}\,\mathord{.}\,\mathord{.})}{\mathrm{d}y}.{}\\
\varphi x'.\varphi x''\mathord{.}\,\mathord{.}\,\mathord{.}\,\mathord{.}\,\mathrm{d}y.\mathrm{d}x'.\mathrm{d}x''\mathord{.}\,\mathord{.}\,\mathord{.}\,\mathord{.}
\end{aligned}
\]
ubi \(y\) extendi debet ab \(y=0\) usque ad \(y=\eta\), indeterminatae reliquae vero per omnes valores, quibus respondet valor realis ipsius \(f(y,x',x''\mathord{.}\,\mathord{.}\,\mathord{.}\,\mathord{.})\). Hinc \SourceLeaf{018}colligitur
\[
\begin{aligned}
\psi y=\int\varphi.f(y,x',x''\mathord{.}\,\mathord{.}\,\mathord{.}\,\mathord{.}).\frac{\mathrm{d}f(y,x',x''\mathord{.}\,\mathord{.}\,\mathord{.}\,\mathord{.})}{\mathrm{d}y}.{}\\
\varphi x'.\varphi x''\mathord{.}\,\mathord{.}\,\mathord{.}\,\mathord{.}\,\mathrm{d}x'.\mathrm{d}x''\mathord{.}\,\mathord{.}\,\mathord{.}\,\mathord{.}
\end{aligned}
\]
integratione, in qua \(y\) tamquam constans considerari debet, extensa per omnes valores indeterminatarum \(x',x''\) etc., qui ipsi \(f(y,x',x''\mathord{.}\,\mathord{.}\,\mathord{.}\,\mathord{.})\) valorem realem conciliant.

\Article{13}
Ad hanc integrationem reipsa exsequendam cognitio functionis \(\varphi\) requireretur, quae plerumque incognita est: quin adeo, etiamsi haec functio cognita esset, in plerisque casibus integratio vires analyseos superaret. Quae quum ita sint, probabilitatem quidem singulorum valorum ipsius \(y\) assignare non poterimus: at secus res se habebit, si tantummodo desideratur valor medius ipsius \(y\), qui oritur ex integratione \(\int y\psi y.\mathrm{d}y\) per omnes valores ipsius \(y\), quos quidem assequi potest, extensa. Et quum manifesto pro omnibus valoribus, quos \(y\) assequi nequit, vel per naturam functionis, quam exprimit (e. g. pro negativis, si esset \(y=xx+x'x'+x''x''+\text{etc.}\)), vel ideo, quod erroribus ipsis \(x,x',x''\) etc. certi limites sunt positi, statuere oporteat \(\psi y=0\), manifesto res perinde se habebit, si integratio illa extendatur per omnes valores reales ipsius \(y\), puta ab \(y=-\infty\) usque ad \(y=+\infty\). Iam integrale \(\int y\psi y.\mathrm{d}y\) inter limites determinatos, puta ab \(y=\eta\) usque ad \(y=\eta'\) sumtum aequale est integrali
\[
\begin{aligned}
\int y\varphi.f(y,x',x''\mathord{.}\,\mathord{.}\,\mathord{.}\,\mathord{.})\frac{\mathrm{d}f(y,x',x''\mathord{.}\,\mathord{.}\,\mathord{.}\,\mathord{.})}{\mathrm{d}y}.{}\\
\varphi x'.\varphi x''\mathord{.}\,\mathord{.}\,\mathord{.}\,\mathord{.}\,\mathrm{d}y.\mathrm{d}x'.\mathrm{d}x''\mathord{.}\,\mathord{.}\,\mathord{.}\,\mathord{.}
\end{aligned}
\]
integratione extensa ab \(y=\eta\) usque ad \(y=\eta'\), atque per omnes valores indeterminatarum \(x',x''\) etc., quibus respondet valor realis ipsius \(f(y,x',x''\mathord{.}\,\mathord{.}\,\mathord{.}\,\mathord{.})\), sive quod idem est, valori integralis
\[
\int y\varphi x.\varphi x'.\varphi x''\mathord{.}\,\mathord{.}\,\mathord{.}\,\mathord{.}\,\mathrm{d}x.\mathrm{d}x'.\mathrm{d}x''\mathord{.}\,\mathord{.}\,\mathord{.}\,\mathord{.}
\]
adhibendo in hac integratione pro \(y\) eius valorem per \(x,x',x''\) etc. expressum, extendendoque eam per omnes harum indeterminatarum valores, quibus respondet valor ipsius \(y\) inter \(\eta\) et \(\eta'\) situs. Hinc colligimus, integrale \(\int y\psi y.\mathrm{d}y\) per omnes valores ipsius \(y\), ab \(y=-\infty\) usque ad \(y=+\infty\) extensum obtineri ex integratione
\[
\int y\varphi x.\varphi x'.\varphi x''\mathord{.}\,\mathord{.}\,\mathord{.}\,\mathord{.}\,\mathrm{d}x.\mathrm{d}x'.\mathrm{d}x''\mathord{.}\,\mathord{.}\,\mathord{.}\,\mathord{.}
\]
